\section{Optimal pair tests with internal role switching}\label{sec:pair-tests}
We now use the full signature $\mathbf A_I^{\rm sym}$, retaining both
winner constants and internal role permutations.

\subsection{Every equation has a singleton-or-pair countervaluation}

Assign each input variable a status: a specified singleton or an
unspecified ambiguous set. Fixing the statuses fixes every test for
whether an intermediate result is a singleton. The remaining dependence
has a restricted positive form.

\begin{lemma}[Permutation provenance on a status cell]\label{sym:lem:provenance}
On every fixed status cell, a term denotes
\begin{equation}\label{sym:eq:provenance}
 C\ \cup\ \bigcup_j\ \bigcup_{\pi\in P_j}\pi S_j,
\end{equation}
where $C\subseteq I$, only ambiguous input variables occur in the union,
and every $P_j$ is a finite, possibly empty, set of permutations.
\end{lemma}
\begin{proof}
A winner or singleton-status input contributes a constant. An ambiguous
input contributes its identity permutation. A role switch transports
$C$ and composes the permutations. A nonempty $P_j$ contributes an
ambiguous set, so that expression cannot equal an owner singleton. Thus
an owned choice collapses exactly when some child is a constant owner
singleton on the cell. Otherwise union its constants and coefficients.
This proves the representation by induction. It is not asserted to be
canonical: distinct coefficients may have the same action on ambiguous
inputs.
\end{proof}

\begin{theorem}[Pair countervaluations]\label{sym:thm:pairs}
If an equation between terms in this full signature fails, it fails at a
valuation in which every variable has one or two possible winners.
\end{theorem}
\begin{proof}
Take a failed valuation, fix its statuses, and choose a role $k$ present
in the left output and absent from the right, interchanging sides if
necessary. Apply Lemma~\ref{sym:lem:provenance} to both terms.

If $k$ lies in the left constant, it survives every valuation in the
cell. Otherwise it has a witness $k=\pi(i)$ for one ambiguous input
$S_j$, one coefficient permutation, and some $i\in S_j$. Replace that
input by a pair contained in it and containing $i$. Replace every other
ambiguous input by any pair contained in it. Leave singleton inputs
fixed. The left witness remains. The right output can only shrink under
these replacements, since every coefficient is a permutation image and
the status cell is unchanged. It still excludes $k$.
\end{proof}

\begin{lemma}[Resetting atomic inputs]\label{sym:lem:reset}
If a unary term function $F$ is extensive, start with $U(x)=F(x)$ and
successively replace $U(x)$ by $U(x)\ow i x$, for each $i\in I$. This gives a term
$\operatorname{Reset}(F)$ which fixes all singletons and agrees with $F$
on every ambiguous input.
\end{lemma}
\begin{proof}
At input $\{i\}$, its own suppression returns $\{i\}$. A different
owner preserves the current value because it contains $i$. On ambiguous
input $S$, the current value contains $S$; neither operand is a
singleton, and their union is still the current value. Extensivity is
preserved at each step.
\end{proof}

\begin{theorem}[Unique minimum fixed test domain]\label{sym:thm:optimal}
A fixed input domain tests every term equation, at every arity, exactly
when it includes every singleton and every two-element set. Consequently
the unique minimum domain is
\[
 D_2=\{S\subseteq I:1\leq|S|\leq2\},\qquad |D_2|=K(K+1)/2.
\]
\end{theorem}
\begin{proof}
Sufficiency is Theorem~\ref{sym:thm:pairs}. Each singleton is necessary
already in the reduct: its unary synthesis gives a term differing from
the constant $I$ only at that singleton. These separators are now
kernel-checked in \texttt{Synthesis.lean}.

For necessity of a particular pair $P$, choose $k\notin P$, put
$C=I\setminus\{k\}$, and form the extensive term
\[
 F(x)=x\cup C\cup\bigcup_{i\in I\setminus(P\cup\{k\})}(i\ k)(x).
\]
Here $(i\ k)$ is a transposition. If $x\subseteq P$, then $F(x)=C$.
Otherwise $x$ contains $k$, or contains an $i$ whose transposition
supplies $k$, and $F(x)=I$. Thus, on ambiguous inputs, $F$ differs from
$I$ only at $P$. By Lemma~\ref{sym:lem:reset},
$\operatorname{Reset}(F)$ and $\operatorname{Reset}(I)$ agree at every
singleton and every ambiguous input other than $P$, and disagree at
$P$. Every complete input domain must therefore contain $P$.
\end{proof}

All $K$ owners and role switches remain available in these tests.
Intermediate outputs may contain any number of winners. The domain
$D_2$ is not closed under the operations; it is an input test domain,
not a replacement algebra. For four roles it has ten values, while the
reduct admits six. The pair separator above explicitly refutes reusing
that six-value test domain after adding internal role switches.

The theorem supplies a decision procedure, uniform in $K$, for term
equality: test $D_2^n$. It does not by itself supply a finite equational
basis or a new proof calculus.
